Source : Eurostat, comptes annuels des administrations publiques (gov\_10a\_main, gov\_10a\_taxag, gov\_10a\_exp), PIB (nama\_10\_gdp) ; Insee, dette trimestrielle de Maastricht par sous-secteur (contributions consolidées, 4e trimestre). Calculs de l'auteur (dette/dynamique.py).\par Lecture : chaque barre est le déficit primaire de l'année (un excédent de 1998 à 2001), réparti selon la part de chaque fonction dans la dépense hors intérêts. Cumulés sur 1996-2025, ces déficits font 50,5 points de PIB, dont 29,7 pour la santé et la protection sociale. Par construction, la répartition de chaque barre suit la structure de la dépense de l'année.
